\customchapter{INTRODUCCIÓN}

La segmentación automática de estructuras óseas en tomografías computarizadas (TC) pélvicas constituye una tarea relevante para el análisis de lesiones, la reconstrucción tridimensional y la planificación de procedimientos quirúrgicos. Los métodos basados en aprendizaje profundo han permitido automatizar esta tarea; sin embargo, su desempeño depende de que los datos de entrenamiento representen las variaciones que pueden encontrarse en las imágenes reales. Entre estas variaciones se encuentran las fracturas, las diferencias entre protocolos de adquisición y la presencia de implantes metálicos, siendo esta última una de las condiciones más difíciles para la segmentación de los huesos pélvicos \cite{liu2021ctpelvic1k}.

Los implantes metálicos, como tornillos, placas y prótesis, presentan una elevada atenuación de los rayos X y producen artefactos que suelen manifestarse como bandas oscuras, regiones brillantes y trazos en forma de rayas. Estos artefactos inducidos por metal pueden ocultar los límites de las estructuras anatómicas y reducir la calidad de las imágenes utilizadas por los modelos de segmentación \cite{selles2024marreview,xie2024implantsegmentation}. Su formación está relacionada con fenómenos físicos y computacionales, entre ellos el endurecimiento del haz, la escasez de fotones y las inconsistencias introducidas durante la reconstrucción. Por esta razón, reproducir de manera exacta el proceso de formación de los artefactos requiere normalmente información del dominio de proyección, como los sinogramas y los parámetros de adquisición.

La disponibilidad de datos anotados bajo esta condición es limitada. CTPelvic1K reúne 1\,184 volúmenes de TC pélvica procedentes de diferentes fuentes, de los cuales 75 contienen artefactos metálicos. El conjunto proporciona anotaciones óseas para 1\,109 volúmenes sin metal y para únicamente 14 volúmenes afectados por metal, mientras que los 61 casos restantes con metal se mantienen sin anotaciones óseas y no se utilizan en el entrenamiento supervisado original \cite{liu2021ctpelvic1k}. Esta distribución produce una representación reducida de los casos metálicos durante el entrenamiento y limita el desarrollo de modelos robustos ante este tipo de alteración.

El aumento de datos mediante transformaciones geométricas o fotométricas convencionales permite incrementar la variabilidad de un conjunto de entrenamiento, pero no modela explícitamente la presencia, la forma ni la ubicación de los implantes metálicos. Tampoco reproduce los patrones de intensidad asociados con los artefactos que estos producen. En este contexto, los modelos de difusión ofrecen un mecanismo para aprender distribuciones complejas y sintetizar imágenes médicas \cite{kazerouni2023diffusionsurvey}. Los modelos de difusión latente reducen el costo computacional al realizar el proceso de difusión en un espacio latente \cite{rombach2022latentdiffusion}, mientras que ControlNet permite incorporar condiciones espaciales, como mapas de segmentación o máscaras, en un modelo de difusión previamente entrenado \cite{zhang2023controlnet}.

La síntesis condicionada también permite evaluar la utilidad de las imágenes generadas mediante una tarea posterior (\textit{downstream task}). En síntesis de tumores, por ejemplo, se ha evaluado si los casos artificiales contribuyen al entrenamiento de modelos capaces de detectar o segmentar lesiones reales \cite{chen2024tumorsynthesis}. Esta forma de evaluación resulta pertinente porque una imagen sintética puede presentar similitud visual sin contener características útiles para un modelo de análisis médico. Esa evaluación presupone, sin embargo, que lo sintetizado sea plausible, y este trabajo examina esa plausibilidad antes que la utilidad (véase la justificación).

A partir de esta necesidad, la presente investigación propone una cadena de síntesis de dos componentes que inserta un tornillo de osteosíntesis y su artefacto metálico local en volúmenes de TC pélvica sin implante. El primer componente, un muestreador de colocación, propone la pose tridimensional del tornillo dentro del corredor óseo medido en cada volumen. El segundo, un sintetizador por difusión, genera su apariencia en el dominio de imagen, sin espacio latente ni datos de proyección. La generación se restringe a la máscara del implante y a una banda a su alrededor, para que el artefacto pueda aparecer fuera del metal, y el resto del volumen se copia sin cambios.

La evaluación mide la coherencia física y quirúrgica de lo sintetizado, y no su utilidad para entrenar una red de segmentación. La coherencia de la colocación se cuantifica comparando la distribución de grados de brecha cortical de las poses, es decir, de cuánto sobresale el implante del hueso, con distribuciones clínicas publicadas. La coherencia de la apariencia se cuantifica con las métricas de un protocolo publicado de simulación física, frente a ese protocolo y frente a una inserción por copia y pegado, que coloca el metal sin artefacto. Antes del sintetizador, una compuerta examina en qué representación de los valores de TC, expresados en unidades Hounsfield (HU), puede operar la síntesis.

El documento presenta, después de este capítulo, los fundamentos teóricos y el estado del arte relacionados con la tomografía, los artefactos inducidos por metal, la fijación iliosacra y los modelos de difusión condicionada. Posteriormente, se describe la metodología propuesta, el diseño experimental y las métricas de evaluación. Finalmente, se presentan y discuten los resultados, seguidos de las conclusiones, limitaciones y posibles líneas de trabajo futuro.

\introsection{Formulación del problema}

¿En qué medida la cadena propuesta genera colocaciones coherentes con las distribuciones clínicas de brecha cortical y artefactos coherentes con el protocolo físico de Peters et al.~\cite{peters2025hybrid}, frente a la inserción por copia y pegado?

\introsection{Objetivos de investigación}

% Fuentes de esta subseccion y de las dos siguientes: tesis/main.tex (Objectives, Out of scope,
% Recognized Gap, Datasets), docs/00-tesis.md, docs/01-decisiones.md y capitulo3.tex.
% Fuente de cada cifra: redaccion/rondas/introduccion-r00-respuesta.md e introduccion-r01-respuesta.md.
Los objetivos son los cuatro que desarrolla el Capítulo~\ref{cap:propuesta}, pero solo los Objetivos~2 y~3 son componentes de la cadena propuesta. El Objetivo~1 es una prueba previa a la cadena. Los Objetivos~1 a~3 tienen en ese capítulo una variable dependiente que los verifica; el Objetivo~4 no tiene variable dependiente propia, porque define o adopta las métricas con que se evalúan los Objetivos~2 y~3.

\paragraph{Objetivo general}

Diseñar, implementar y evaluar una cadena de síntesis por difusión en el dominio de imagen que genere implantes metálicos de osteosíntesis y su artefacto metálico local en volúmenes de TC pélvica, con coherencia física y quirúrgica. Los implantes de osteosíntesis son dispositivos de fijación ósea, como tornillos y placas. La colocación del implante se restringe con el corredor óseo medido sobre cada volumen, es decir, la región de hueso por la que puede pasar un tornillo sin atravesar la cortical. La apariencia se condiciona con una codificación multiventana de los HU, formada por varios canales, cada uno con una ventana de HU distinta. El único implante que se sintetiza es un tornillo transilíaco-transsacro, o transiliosacro (véase el alcance).

\paragraph{Objetivos específicos}

\begin{enumerate}
    \item \textbf{Evaluar si el autoencoder de un modelo de difusión latente conserva los HU de una codificación multiventana (compuerta \emph{Go/No-Go}).} Determinar si el paso de ida y vuelta (HU $\rightarrow$ representación $\rightarrow$ HU) por ese autoencoder conserva los HU del hueso con un error absoluto medio menor que 25~HU. El autoencoder es la red que comprime la imagen a un espacio latente, donde opera la difusión, y la reconstruye desde él \cite{rombach2022latentdiffusion}. La regla de la compuerta examina seis combinaciones, de dos variantes del autoencoder y tres codificaciones, sobre 34 pacientes de prueba, y aprueba si alguna cumple el criterio (Sección~\ref{sec:obj1}).

    \item \textbf{Diseñar un muestreador de colocación quirúrgicamente restringido.} Proponer poses tridimensionales de un tornillo transilíaco-transsacro, que cruza las dos articulaciones sacroilíacas de un ilion al otro y se modela como un cilindro paramétrico y rígido, alrededor del corredor óseo medido en cada volumen. Las poses se expresan en el marco de referencia de Kaiser et al.~\cite{kaiser2014dysmorphism}, y un corredor es viable si admite el calibre del tornillo más una holgura radial tomada de esa misma fuente. Cada pose recibe un \textbf{grado de brecha cortical}, es decir, su nivel en una escala ordinal de cuatro niveles según cuánto sobresale el implante del hueso. La distribución de grados se compara con la \textbf{referencia clínica}, formada por las series con navegación computarizada (navegada) y sin ella (convencional) de Zwingmann et al.~\cite{zwingmann2009navigated} en S1, el primer segmento del sacro. La comparación usa la distancia de Wasserstein-1, que sobre grados ordinales suma las diferencias absolutas entre las dos distribuciones acumuladas.

    \item \textbf{Diseñar un sintetizador condicionado.} Formular la síntesis como \emph{inpainting} por difusión 2.5D en el dominio de imagen, es decir, sobre los valores de TC por vóxel, sin espacio latente ni datos de proyección. El \emph{inpainting} regenera solo una región enmascarada de la imagen y conserva el resto; el término 2.5D se precisa en el marco teórico. Lo que se genera es la \textbf{región de generación} $G = M \cup B_{\delta}$, donde $M$ es la \textbf{máscara del implante} y $B_{\delta}$ una \textbf{banda de generación extendida} (\emph{generation band}) de unos 12~mm alrededor de ella. La banda permite que las rayas (\emph{streaking}) y las zonas oscuras del endurecimiento del haz aparezcan fuera del metal. La apariencia generada se compara con la \textbf{inserción por copia y pegado}, que coloca el metal sin artefacto, con el protocolo físico de Peters et al.~\cite{peters2025hybrid} y con observaciones reales de artefacto; en estas últimas, con estadísticos de los HU alrededor del metal que no exigen una imagen sin metal del mismo paciente (Sección~\ref{sec:apariencia}).

    \item \textbf{Formalizar la métrica de colocación y adoptar las métricas de apariencia de Peters et al.} Formalizar la admisibilidad quirúrgica de la colocación (SAP, \emph{Surgical Admissibility of Placement}), la única métrica que introduce este trabajo. SAP reúne dos componentes, definidos en el Objetivo~2: el grado de brecha cortical de cada pose y la viabilidad del corredor, que se evalúa con un calibre de 7.0~mm y una holgura radial de 1~mm (Sección~\ref{sec:sap}). Solo la distribución de grados se compara con la referencia clínica; la viabilidad se reporta de forma descriptiva. La definición de SAP incluye la distancia de Wasserstein-1 de esa comparación, que se ejecuta en el Objetivo~2. La apariencia se evalúa con las métricas de Peters et al., con sus nombres publicados: \emph{bone integrity}, \emph{metal integrity} y \emph{streak amplitude}. Esa adopción no se presenta como contribución evaluada. Esas métricas se diseñaron para evaluar la reducción de artefactos metálicos (MAR, \emph{metal artifact reduction}), y usarlas para evaluar síntesis exige invertirlas. La inversión está definida para la \emph{streak amplitude}; la \emph{bone integrity} y la \emph{metal integrity} conservan su fórmula e invierten su lectura, cuya referencia es el valor del protocolo físico (Sección~\ref{sec:apariencia}).
\end{enumerate}

La compuerta del Objetivo~1 se formuló para una síntesis por difusión latente guiada con ControlNet. Su veredicto fue negativo, y después de conocerlo el Objetivo~3 se reformuló en el dominio de imagen, sin modificar la regla de la compuerta. Con la misma regla y sobre los mismos pacientes de prueba se examinó después un autoencoder entrenado con TC (Sección~\ref{sec:obj1}). El veredicto descartó la ruta latente, no la codificación multiventana, que el sintetizador sigue usando, ahora sin la etapa del autoencoder. La ida y vuelta de esa codificación sin autoencoder se midió en los experimentos del Objetivo~1, aunque no entra en su criterio de fallo (Sección~\ref{sec:obj1}). Sobre los 34 pacientes de prueba, la media del error absoluto medio en el hueso fue de 0.00~HU en dos de las tres codificaciones y de 21.80~HU en la tercera. El sintetizador adopta una de esas dos, la de compresión arcoseno hiperbólico, en coma flotante de 32 bits (Sección~\ref{sec:sintetizador}).

No todos los objetivos fijan de antemano qué resultado los haría fallar; el Objetivo~1 sí lo fija. Su regla se escribió después de una exploración que incluía a los pacientes de prueba y cuyo error ya superaba el criterio, pero antes de la prueba que decide. El criterio de 25~HU es anterior a esa exploración y no se movió (Sección~\ref{sec:obj1}). El Objetivo~2 compara su distribución de grados con la referencia clínica sin regla de decisión, que fijada ahora sería posterior al resultado, y la lee contra la distancia entre las dos series de esa referencia (Sección~\ref{sec:sap}). En el Objetivo~3 decide la comparación de equivalencia frente al protocolo físico: solo se concluye equivalencia si el intervalo de confianza del 90\,\% de la diferencia pareada cae entero en un margen $[-\Delta, +\Delta]$, aún sin medir (Sección~\ref{sec:apariencia}). La superioridad frente a la inserción por copia y pegado, que no produce rayas por construcción, es un control de cordura que supera cualquier brazo que genere algo, y ningún resultado suyo cuenta a favor del sintetizador. El Objetivo~4 es un objetivo de definición: se verifica con los controles de implementación de SAP (identidad, monotonía y resolución), se ejerce en el Objetivo~2 y no se evalúa experimentalmente por separado.

\introsection{Justificación}

Este trabajo se justifica por una carencia de datos anotados. La segmentación ósea automática de la TC pélvica depende de volúmenes anotados, y los artefactos metálicos pueden ocultar los límites anatómicos cerca de los implantes. En CTPelvic1K, el conjunto que usa este trabajo, Liu et al.~\cite{liu2021ctpelvic1k} publican anotaciones óseas para 14 de los 75 volúmenes con artefacto metálico y dejan sin anotar los 61 restantes. La motivación de fondo es usar implantes sintéticos como aumento de datos para la segmentación ósea alrededor del implante, pero esa utilidad no se mide aquí (véase el alcance). Este trabajo evalúa, en su lugar, si lo sintetizado es coherente con la física del artefacto y con la práctica quirúrgica, porque medir la utilidad para segmentación de ejemplos cuya plausibilidad física no se conoce confundiría las dos propiedades.

La justificación computacional es que ninguno de los tres enfoques revisados (simulación física, planificación quirúrgica determinista y síntesis de lesiones por difusión) produce a la vez una distribución de poses del implante y la apariencia de su artefacto. Ninguno de los dos primeros produce esa distribución de poses. En la simulación física, el protocolo físico de Peters et al.~\cite{peters2025hybrid} simula la formación del artefacto, pero coloca el metal al azar en sus datos de entrenamiento, porque la colocación manual les resultaba impracticable. En su evaluación, sobre conjuntos de pacientes que el artículo llama clínicamente representativos, el metal se inserta en ubicaciones que llama pertinentes (\emph{meaningful locations}), sin que describa la regla con que se eligen. Los planificadores quirúrgicos deterministas, como el de Liu et al.~\cite{liu2025pipeline}, optimizan una única trayectoria. Zwingmann et al.~\cite{zwingmann2009navigated}, en cambio, reportan distribuciones de malposición distintas según la técnica quirúrgica, y este trabajo lee ese resultado como indicación de que el muestreador ha de reproducir una distribución de poses y no una pose única.

Este trabajo prevé ejecutar, además, el protocolo de Peters et al. sobre las poses del muestreador, en los pacientes de prueba sin metal y según el plazo, como comparación para la apariencia (Sección~\ref{sec:apariencia}). Ese protocolo no sustituye al sintetizador: necesita simular proyecciones para cada pose, simula en dos dimensiones una sola fila de detector y los objetos metálicos de sus datos de entrenamiento son formas fractales \cite{peters2025hybrid}. El sintetizador, en cambio, opera sobre cualquier TC reconstruida, sin proyecciones. Si alcanza la equivalencia con la simulación física es justamente lo que pone a prueba el Objetivo~3, de modo que esa equivalencia es una pregunta de este trabajo y no una premisa.

El tercer enfoque, la síntesis de lesiones por difusión, no liga al objeto generado los cambios de intensidad fuera de su máscara, y el artefacto metálico no queda dentro del implante. De Man et al.~\cite{deman1999} identifican el endurecimiento del haz, la dispersión, el ruido y el efecto exponencial de gradiente de borde como las causas principales de las rayas del metal, y las muestran irradiando desde el metal. Los modelos de difusión revisados generan la lesión por \emph{inpainting} acotado a su máscara \cite{ramzan2026claim,chen2024tumorsynthesis,jacob2026lgesynthnet} o sintetizan el corte entero desde ruido \cite{zhang2025diffboost}. En los que acotan el \emph{inpainting} a la máscara no hay un mecanismo para los cambios de intensidad que el objeto generado induce fuera de ella, y en ninguno hay una señal de entrenamiento para esos cambios.

La contribución de este trabajo es la combinación de tres elementos que las fuentes revisadas no reúnen. El primero son geometrías de implante rígidas y paramétricas como máscara de síntesis, que evitan segmentar con un umbral fijo de HU el metal que se coloca. Xie et al.~\cite{xie2024implantsegmentation} reportan que esa segmentación, en cortes simulados, cubrió el implante en exceso, y en la cohorte local los objetos metálicos alargados extraídos con un umbral fijo aparecen fragmentados (Sección~\ref{sec:datos}). El entrenamiento del sintetizador sí usa máscaras umbralizadas de metal real, y su diferencia con el cilindro paramétrico es uno de los desplazamientos de dominio entre entrenamiento y síntesis (Sección~\ref{sec:sintetizador}). Ningún objetivo aísla este aporte, porque ningún brazo de comparación coloca geometrías extraídas por umbral (Sección~\ref{sec:geometria}).

El segundo elemento es un muestreador cuya distribución de poses se compara con la referencia clínica, y el tercero, la codificación multiventana usada para generar, junto con la banda $B_{\delta}$. En las fuentes revisadas no hay un precedente que fije con un valor numérico una banda de generación fuera de la máscara del implante. El uso de varias ventanas de HU proviene de trabajos de MAR, que las aplican a etapas de reconstrucción o a la función de pérdida y no a la codificación de la entrada \cite{wang2025adaptiveweighting,li2024}. El aporte de este trabajo se limita a emplear esas ventanas como codificación de la entrada para generar. Del sintetizador se reclama la combinación del objeto metálico rígido, la codificación multiventana, la banda $B_{\delta}$ y la copia exacta del volumen fuera de $G$, y no cada pieza por separado (Sección~\ref{sec:ea-brecha}). Ningún objetivo aísla la codificación ni la banda, porque el Objetivo~3 evalúa el sintetizador completo.

La justificación metodológica tiene dos partes: el muestreador no se construye con los datos contra los que se compara, y el veredicto negativo del Objetivo~1 aporta un hallazgo propio. La distribución de poses del muestreador se fijó por escrito en su preinscripción, antes de calcular cualquier distancia contra la referencia clínica, y ningún parámetro del muestreador se toma de ella. El hallazgo del Objetivo~1 es que ninguna de las seis combinaciones bajó de 25~HU, y que en el autoencoder entrenado con TC la sola saturación a su rango de salida ya superaba ese criterio (Sección~\ref{sec:obj1}).

\introsection{Alcance y limitaciones / restricciones}

El alcance queda acotado en cuatro puntos: la cohorte, el implante, el nivel sacro y la forma de la síntesis. La anatomía receptora y las observaciones reales de artefacto provienen del subconjunto local de CTPelvic1K, 178 de los 1\,184 volúmenes que declara la publicación \cite{liu2021ctpelvic1k}. El único implante que se coloca y se sintetiza es un tornillo transilíaco-transsacro paramétrico y rígido; los implantes reales del conjunto no aportan geometría. Es el único implante con un corredor óseo medible sobre cada volumen y con una distribución clínica publicada de grados de brecha, y en las fuentes revisadas las placas y los demás tornillos no tienen ninguna de las dos. Esa distribución, como el marco de Kaiser et al., se refiere a tornillos iliosacros, y entra como referencia externa de grados y no como prueba de que ambos tornillos sean el mismo (Sección~\ref{sec:mt-iliosacra}).

En el nivel sacro, la comparación con la referencia clínica se limita a S1. El segundo corredor bajo S1, un corredor óseo cuyo nivel vertebral varía entre volúmenes, queda fuera de esa comparación, porque ninguna de las fuentes revisadas reporta para él una distribución ordinal clínica. Su eje está medido, pero el muestreo de poses en él no se ha ejecutado \GAPDATO{preinscripción y corrida del muestreo de poses en el segundo corredor bajo S1}. La síntesis opera sobre parches alrededor del implante, y dentro de cada parche solo genera la región $G$; fuera de ella copia el volumen de origen (Sección~\ref{sec:sintetizador}).

Quedan fuera de alcance cuatro tareas, cada una por una razón distinta:
\begin{itemize}
    \item \textbf{Evaluación de segmentación posterior.} Medir el efecto de los implantes sintéticos sobre una red de segmentación ósea con el coeficiente de Dice y la distancia de Hausdorff al percentil~95 (HD95) exigiría casos reales con metal y anotación ósea. En CTPelvic1K son 14 (véase la justificación), y no está verificado que sus anotaciones estén disponibles localmente. Además, localmente solo se dispone de una parte de la colección. Un resultado de segmentación descansaría en una cohorte pequeña y parcialmente anotada.
    \item \textbf{Ablaciones del muestreador.} Retirar uno a uno los componentes del muestreador vigente, el anclaje al eje del corredor y la escala de la perturbación, para medir el efecto de cada uno se difiere a trabajo futuro, para concentrar el tiempo en la compuerta y el muestreador.
    \item \textbf{Reimplementación validada de XCIST/CatSim.} La validación publicada del simulador es preliminar y no incluye un estudio de artefacto metálico \cite{wu2022xcist}, de modo que una reimplementación propia no tendría ninguna validación en metal. En su lugar se adopta el protocolo físico de Peters et al.~\cite{peters2025hybrid}, que se ejecuta sobre ese mismo simulador, publica sus métricas y valida su simulación de artefactos metálicos contra un fantoma físico. Esa validación tampoco cubre la síntesis (véanse las limitaciones).
    \item \textbf{Estratificación por dismorfismo sacro.} Gardner et al.~\cite{gardner2010safezones} clasifican el sacro como dismórfico sobre radiografía simple, por rasgos como el ala sacra angulada y ascendente o los forámenes neurales superiores no circulares, sin exigir que estén todos presentes. El efecto del dismorfismo sobre el tamaño de la zona segura no es consistente entre estudios. Esos autores reportan en S1 un área transversal mínima de la zona segura de 222~mm$^2$ en sacros dismórficos frente a 346~mm$^2$ en normales, pero recogen un estudio previo que no halló diferencia. Por eso el muestreador no estratifica por dismorfismo y mide, en cambio, la geometría del corredor sobre cada volumen.
\end{itemize}

El diseño descansa en supuestos, que son afirmaciones no verificadas sobre la física del artefacto o sobre las fuentes, y en convenciones, que son valores fijados por este trabajo; la Sección~\ref{sec:amenazas} discute su efecto sobre los resultados. El primer supuesto es que la apariencia del artefacto metálico puede generarse en el dominio de imagen, sin datos de proyección, aunque De Man et al.~\cite{deman1999} caracterizan sus mecanismos como no lineales y los aíslan modificando el sinograma. La comparación con el protocolo físico es la que puede contradecirlo (véanse las limitaciones). El segundo supuesto es que la serie navegada de la referencia clínica, cuyo nivel sacro Zwingmann et al. no declaran, corresponde a S1. El tercero es que la relación entre la máscara y el artefacto no depende del tipo de implante, lo que permite entrenar el sintetizador con componentes metálicos de todo tipo.

Además de los valores declarados en la preinscripción del muestreador (Tabla~\ref{tab:preinscripcion}), el diseño fija otras convenciones. Ninguna fuente revisada calibra el ancho de unos 12~mm de $B_{\delta}$. Los límites de grado de la escala de brecha cortical, de 2~mm de ancho, provienen de la literatura de tornillos pediculares \cite{gertzbein1990,mirza2003} y llegan a la fijación iliosacra a través de Smith et al.~\cite{smith2006iliosacral}. Su traslado al corredor iliosacro es una convención geométrica y un supuesto explícito, no un umbral clínico validado. Los cuatro grados de brecha se tratan como equiespaciados al calcular la distancia de Wasserstein-1. La desviación estándar de las perturbaciones del eje se fija en la mitad del margen cortical de longitud útil de 5~mm que exigen Kaiser et al.~\cite{kaiser2014dysmorphism}. Esa convención no tiene fuente y se justifica por su consecuencia: con perturbaciones de magnitud seminormal, cerca del 95\,\% de las magnitudes quedan dentro de ese margen (Sección~\ref{sec:poses}).

Las limitaciones se resumen aquí por componente, y la Sección~\ref{sec:amenazas} las discute en detalle. En la colocación, la referencia clínica proviene de pelvis fracturadas, mientras que las pelvis receptoras se eligieron sin osteosíntesis y no se verificaron libres de fractura. Una malreducción residual en la referencia clínica estrecharía los corredores de sus pacientes de una forma que las pelvis receptoras no reproducen, y una fractura no detectada en estas los estrecharía por una causa distinta de la variabilidad anatómica. En un corredor más estrecho que el tornillo con que se mide la brecha, el propio eje ya perfora y el grado~0 es inalcanzable por anatomía (Sección~\ref{sec:sap}). Por eso ambas diferencias pueden mover la distancia de Wasserstein-1 sin que intervenga el muestreador. Un cribado ciego de fractura sobre 30 casos de la cohorte del Objetivo~2, 15 de corredor estrecho y 15 de control, halló fractura confirmada en 20 de los 30, 10 en cada grupo (Sección~\ref{sec:datos}). Lo hizo un médico sin especialidad, de modo que es un cribado y no una lectura de especialista, y muestra que la cohorte receptora no es una colección de pelvis sanas.

Las cifras principales del Objetivo~2 dependen de una variante de la segmentación anatómica cuya elección quedó condicionada a revisar 16 casos, marcados por diferencia de diámetro entre variantes o por desplazamiento de la caja de una estructura. La autora, con apoyo de un médico sin especialidad, revisó los 16 sobre láminas: 14 salieron conformes y 2 fallaron con las dos variantes (Sección~\ref{sec:corredor}). Ninguno falló solo con la variante principal, que es el único veredicto que habría reabierto su elección, y la variante se mantiene. La publicación de TotalSegmentator, la herramienta de esa segmentación, no reporta exactitud para la etiqueta de S1 \cite{wasserthal2023}, de modo que esa exactitud no se asume y un control verifica el nivel S1 de cada caso. Entre los candidatos a la cohorte del Objetivo~2, la discordancia de nivel excluyó a 11 de 34 pacientes con objetos no ortopédicos y a 7 de 57 sin objeto. Si el corredor difiriera entre esos grupos, la distribución de grados heredaría la composición de la cohorte. La brecha se mide contra una envolvente ósea, una máscara de hueso cuyo cierre morfológico podría ocupar el canal sacro y los forámenes, y en ese caso una perforación hacia ellos daría brecha nula \GAPDATO{comprobación de que la envolvente ósea deja fuera el canal sacro y los forámenes sacros}.

En la apariencia, las observaciones reales de artefacto tienen una reconstrucción no reportada, que pudo incluir una MAR del fabricante o una reconstrucción monoenergética, y ambas cambian la apariencia del artefacto \cite{selles2024marreview}. Por eso sirven para comparar la apariencia y no como verdad física. El protocolo físico adoptado se validó en dos dimensiones y para MAR, sin cubrir el paso híbrido con imágenes clínicas, de modo que usarlo como comparación para la síntesis exige una adaptación que no hereda esa validación. Como $B_{\delta}$ trunca por diseño las rayas lejanas y el protocolo físico no, la comparación entre ambos no puede mostrar si el sintetizador reproduciría las rayas más allá de la banda.

Entre el entrenamiento del sintetizador y la síntesis hay tres desplazamientos de dominio, que limitan la transferencia de lo aprendido sobre metal real a los tornillos paramétricos. Las máscaras de entrenamiento se umbralizan sobre metal real, y la de síntesis es el cilindro paramétrico liso. El contexto de entrenamiento trae rayas lejanas que no tienen los volúmenes sin metal que reciben la síntesis. Los parches que solo contienen banda son menos frecuentes en el entrenamiento que en la síntesis. El primero se cuantifica con la fragmentación y el diámetro de los tornillos medidos en la auditoría, el segundo se mide en la evaluación del sintetizador y el tercero, con la razón entre tipos de parche (Sección~\ref{sec:sintetizador}). Por último, el Objetivo~3 no tiene todavía resultados \GAPDATO{el sintetizador generó muestras sobre componentes metálicos reales y la cadena completa se ejecutó sobre pacientes de validación, pero no existe aún resultado del Objetivo~3 porque el modelo final no está elegido}, y su comparación de equivalencia frente al protocolo físico depende del plazo \GAPDEC{si la comparación de equivalencia frente al protocolo físico se mantiene o pasa a trabajo futuro, contingencia de plazo registrada por la autora}.
